\documentclass[11pt]{article}

\usepackage[margin=1in]{geometry}
\usepackage{graphicx}
\usepackage{booktabs}
\usepackage{amsmath}
\usepackage{amssymb}
\usepackage{natbib}
\usepackage{hyperref}
\usepackage{xcolor}
\usepackage{microtype}
\bibliographystyle{unsrtnat}

\title{A SUMO E3 ligase promotes long non-coding RNA transcription to
regulate small RNA-directed DNA elimination}
\author{}
\date{}

\begin{document}
\maketitle

\begin{abstract}
Small RNAs loaded onto Argonaute proteins silence complementary genomic
regions by recognising nascent long non-coding RNAs (lncRNAs) at chromatin,
which creates the apparent paradox that loci destined for silencing must
first be transcribed. In the ciliate \emph{Tetrahymena thermophila},
programmed DNA elimination depends on parental macronuclear lncRNAs
(pMAC-lncRNAs) that base-pair with scan RNAs (scnRNAs) to trigger
target-directed small RNA degradation (TDSD), yet no factor that licenses
pMAC-lncRNA transcription has been identified. Here we characterise Ema2
(TTHERM\_00113330), a conjugation-specific SUMO E3 ligase that carries an
atypical SP-RING domain in which the first zinc-coordinating cysteine is
replaced by histidine. Using somatic knockout strains, in vitro GST
pull-down, mass spectrometry of His-Smt3-conjugated proteins, site-directed
lysine-to-arginine mutagenesis of the transcription elongation factor Spt6,
and biochemical fractionation, we show that Ema2 interacts with the SUMO E2
Ubc9, catalyses SUMOylation of Spt6, and sustains the chromatin association
of Spt6 and RNA polymerase II (RNAPII) during mid-conjugation.
\emph{EMA2} knockout cells fail to complete DNA elimination, lose TDSD of
MDS-complementary scnRNAs, abolish pMAC-lncRNA accumulation and parental-MAC
dsRNA, and produce no viable sexual progeny; at the biochemical level,
high-molecular-weight ($>\!200$~kDa) SUMO conjugates are reduced to
approximately $50\%$ of wild-type levels. A SUMOylation-defective Spt6
(C-KR) mutant preserves pMAC-lncRNA transcription but still shows a mild
DNA elimination defect, demonstrating that Spt6 SUMOylation is dispensable
for lncRNA production yet contributes to later elimination steps. These
findings place Ema2-directed SUMOylation upstream of lncRNA transcription
and define a heterochromatin-independent solution to small RNA-directed
genome communication.
\end{abstract}

\section{Introduction}

Small RNAs complexed with Argonaute family proteins silence complementary
genomic regions either by degrading messenger RNAs in the cytoplasm or by
guiding transcriptional silencing at chromatin
\citep{holoch2015rna,wilson2013rnai}. Transcriptional silencing presents a
conceptual puzzle because the Argonaute--small RNA complex must recognise
its target through a nascent long non-coding RNA (lncRNA), which means
that a region destined for silencing must simultaneously be transcribed
by RNA polymerase. Different lineages have evolved different compromises
with this paradox, but each solution assumes that lncRNA production at the
silenced locus is somehow licensed. Understanding the factors that license
that transcription is therefore central to understanding how small RNAs
communicate with chromatin.

Comparisons across model systems reveal how varied these licensing
strategies can be. In fission yeast, small interfering RNAs (siRNAs)
direct the deposition of histone H3 lysine~9 di- and tri-methylation
(H3K9me2/3) on centromeric repeats, leading to heterochromatin assembly
by the HP1 protein Swi6 \citep{volpe2002rnai,hall2002heterochromatin};
mitotic phosphorylation of histone H3 serine~10 transiently evicts Swi6 so
that RNAPII can re-transcribe the repeats during a narrow cell-cycle
window. In \emph{Drosophila}, the HP1 paralog Rhino binds H3K9me2/3 at
piRNA clusters \citep{lethomas2014rhino} and recruits dedicated variants
of basal transcription factors to sustain lncRNA production from
heterochromatin \citep{andersen2017heterochromatin}, while the SUMO E3
ligase Su(var)2-10 couples piRNA target recognition to SetDB1 recruitment
and reinforces the silenced state \citep{ninova2019suvar}. In plants,
SHH1 delivers the plant-specific RNA polymerase~IV to H3K9-methylated loci
to licence siRNA precursor transcription for RNA-directed DNA methylation
\citep{law2013shh1}. A unifying feature of these systems is that
heterochromatin is already in place when lncRNA transcription is
licensed, so the mechanisms that initiate silencing and those that sustain
it are difficult to separate.

Ciliates circumvent this confound because the source and target of the
silencing small RNA occupy physically distinct nuclei within a single
cell \citep{chalker2011ciliate}. Each \emph{T.~thermophila} cell carries a
diploid germline micronucleus (MIC) and a polyploid somatic macronucleus
(MAC); during conjugation the MIC undergoes meiosis and fertilisation to
produce a new MIC and a new MAC, while the parental MAC is destroyed.
Programmed DNA elimination takes place in the developing new MAC and
excises roughly $12{,}000$ internal elimination sequences (IESs), reducing
the approximately $200$~Mb micronuclear genome to the approximately
$103$~Mb macronuclear genome \citep{cheng2010transposase,dubois2017piggymac}.
Three classes of lncRNAs orchestrate this rearrangement: bidirectional
MIC-lncRNAs produced during meiotic prophase
\citep{mochizuki2004rpb3,chalker2001nongenic} are processed into
approximately $29$~nt scnRNAs loaded onto the Piwi protein Twi1
\citep{mochizuki2002twi1,noto2010giw1}; after Twi1--scnRNA complexes enter
the parental MAC, pMAC-lncRNAs base-pair with them and, via the RNA
helicase Ema1 \citep{aronica2008ema1}, trigger TDSD so that only
IES-matching scnRNAs are retained \citep{schoeberl2012biased}; finally,
nMAC-lncRNAs in the new MAC recruit the Polycomb Repressive Complex~2
homolog Ezl1, which deposits H3K9me2/3 and H3K27me3 to mark IESs for
excision by domesticated PiggyBac transposases
\citep{liu2004h3k9,liu2007rnai,frapporti2019ezl1,cheng2010transposase,dubois2017piggymac}.

Despite the central role of pMAC-lncRNAs in TDSD, no factor that promotes
their transcription has been described. SUMOylation is a strong candidate
regulator because it increases sharply during \emph{Tetrahymena}
conjugation \citep{nasir2015sumo}, because RNAi knockdown of the SUMO
pathway blocks IES excision in the related ciliate \emph{Paramecium
tetraurelia} \citep{matsuda2006sumo}, and because SUMO modifications of
chromatin factors are recurrently associated with small RNA-directed
silencing in fission yeast \citep{shin2005sumo}, \emph{C.~elegans}
\citep{kim2021hdac1} and \emph{Drosophila} \citep{ninova2019suvar}.
The \emph{Tetrahymena} genome encodes a single conjugation-specific SUMO
E3 ligase of the Siz/PIAS family, Ema2 (TTHERM\_00113330), whose SP-RING
signature is atypical \citep{hochstrasser2001spring}. Whether Ema2 is a
bona fide SUMO E3, and whether it acts on the lncRNA pathway rather than
on a downstream excision step, has not been tested.

In this study we take Ema2 as an entry point into the regulation of
pMAC-lncRNA transcription. We combine somatic knockouts, immunofluorescence
and DNA fluorescence in situ hybridisation, northern blotting and small
RNA sequencing, in vitro GST pull-down with recombinant Ubc9, mass
spectrometry of His-Smt3 conjugates under denaturing conditions, cell
fractionation with and without RNase treatment, and site-directed
lysine-to-arginine mutagenesis of Spt6 to establish that Ema2 catalyses
conjugation-stage SUMOylation, targets the transcription elongation factor
Spt6, and sustains chromatin association of Spt6 and RNAPII during
mid-conjugation. We further show that a SUMOylation-defective Spt6
(C-KR) mutant preserves pMAC-lncRNA transcription while still displaying
a mild DNA elimination defect, indicating that Ema2 acts through multiple
SUMO substrates and that Spt6 SUMOylation itself has a distinct, later
role in elimination. Together our results identify a SUMO E3 ligase that
operates upstream of lncRNA transcription and offer a mechanistically
distinct solution to the silencing paradox, one that does not presuppose
a heterochromatic template.

\section{Materials and Methods}

\subsection{Strains and culture}
Wild-type \emph{Tetrahymena thermophila} strains B2086, CU427, CU428,
B*VI and B*VII were obtained from the Tetrahymena stock center. \emph{EMA2}
germline knockout strains were those described in Shehzada \& Mochizuki
(2022). Cells were grown in SPP medium containing 2\% proteose peptone at
30~\textdegree C. For conjugation, growing cells at $\sim5$--$7\times10^5$
cells/mL of two mating types were washed, pre-starved for $\sim12$--$24$~h
and mixed in 10~mM Tris (pH~7.5) at 30~\textdegree C. Time points are
reported as hours post-mixing (hpm). Mating pairs for viability tests were
isolated into drops of 1$\times$SPP at 6~hpm and incubated for
approximately three days, after which sexual progeny were identified by
6-methylpurine resistance and paromomycin sensitivity as in
\citet{mochizuki2002twi1}.

\subsection{Antibodies}
Primary antibodies were: anti-H3K27me3 (Millipore 07-449), anti-HA HA11
(Covance 6B12), anti-H3K9me3 (Active Motif 39162), anti-GST (BD 554805),
anti-His (Proteintech 66005-1-Ig), anti-long-dsRNA J2 (Jena Bioscience),
anti-$\alpha$-tubulin 12G10, anti-Rpb3 (Kataoka \& Mochizuki 2017) and
anti-Smt3 (kindly provided by James Forney, as used in
\citealp{nasir2015sumo}).

\subsection{Small RNA analyses}
For northern hybridisation, 10~\textmu g of total RNA extracted with TRIzol
from conjugating wild-type or \emph{EMA2} knockout cells was resolved on a
15\% acrylamide--urea gel and probed with a radiolabelled Mi-9 probe
\citep{aronica2008ema1}. For small RNA sequencing, 26- to 32-nt reads that
uniquely matched a MDS or IES tile were counted as reads per kilobase of
unique sequence per million mapped reads (RPKM).

\subsection{DNA fluorescence in situ hybridisation and IES retention}
Probes for the Tlr1 locus were generated by nick translation with Cy5
from the pMBR~4C1, pMBR~2 and Tlr~IntB plasmids and hybridised to fixed
exconjugants at 36~hpm. Because DNA elimination is completely blocked in
\emph{TWI1} knockouts, a correction factor of 0.35 was applied so that
the mean IES retention index of \emph{TWI1} knockout cells equals 1; other
strains were scored relative to this baseline.

\subsection{GST pull-down assay}
\emph{E.~coli} codon-optimised synthetic \emph{EMA2} (EMA2-Ec) and
\emph{UBC9} (UBC9-Ec) genes were cloned into pGEX-4T-1 (BamHI/XhoI) and
pET28a($+$) (NdeI/BamHI), respectively. His-Ubc9 expression in BL21 was
induced for $\sim16$~h at 25~\textdegree C with 0.05~mM IPTG; cells were
lysed in HEPES buffer (50~mM HEPES pH~7.5, 800~mM KCl, 10\% glycerol,
0.2~mM PMSF, 1$\times$cOmplete EDTA-free) with 10~mM imidazole, bound to
Ni-NTA for 2~h at 4~\textdegree C, washed with buffer containing 20~mM
imidazole, eluted with 250~mM imidazole, dialysed into 50~mM Tris pH~7.5,
150~mM NaCl, 2~mM MgCl$_2$, 1~mM DTT, 50\% glycerol and concentrated
through a Microcon 10-kDa filter. GST-Ema2 or GST alone was bound to
Glutathione Sepharose 4B for 90~min at 4~\textdegree C and incubated with
10~\textmu g of His-Ubc9 in GST pull-down buffer (20~mM Tris--HCl pH~7.5,
100~mM NaCl, 1\% NP-40, 5~mM MgCl$_2$, 1~mM DTT, 0.2~mM PMSF,
1$\times$cOmplete EDTA-free) for 90~min at 4~\textdegree C; beads were
washed four times and eluted with 2$\times$SDS sample buffer.

\subsection{His-Smt3 purification under denaturing conditions}
Cells expressing His-Smt3 were lysed in Guanidine Lysis Buffer (6~M
guanidine-HCl, 93.2~mM Na$_2$HPO$_4$, 6.8~mM NaH$_2$PO$_4$, 10~mM Tris-HCl
pH~8.0) by probe sonication (Omni-Ruptor~250 Ultrasonic Homogenizer;
$2\times25$~pulses, 20\% power, 60\% pulse). Cleared lysates were
supplemented with 50~mM imidazole and 5~mM $\beta$-mercaptoethanol,
incubated with Ni-NTA agarose, washed and eluted with 6~M urea,
58~mM Na$_2$HPO$_4$, 42~mM NaH$_2$PO$_4$, 10~mM Tris-HCl pH~8.0,
500~mM imidazole. Eluted proteins were precipitated in 10\% TCA,
resuspended in 1$\times$SDS buffer and analysed by Western blotting and
mass spectrometry. For the candidate-SUMO-substrate list, we excluded
proteins with a log2 LFQ score above 25 or more than six consecutive
histidine residues to remove common contaminants.

\subsection{HA-Smt3 purification}
Total proteins were precipitated in 6\% TCA on ice for 10~min and pelleted
at 13{,}000~rpm for 5~min at 4~\textdegree C. Pellets were dissolved in
1$\times$SDS, neutralised with 2~M Tris, and heated 5~min at
95~\textdegree C; $600$~\textmu L of dissolved protein was diluted into
7.4~mL of 20~mM Tris pH~7.5, 100~mM NaCl, 2~mM MgCl$_2$, 2~mM CaCl$_2$
and incubated overnight at 4~\textdegree C with Ezview Red anti-HA
agarose (Sigma). Beads were washed five times with the same buffer
containing 0.1\% Tween-20, then eluted by heating at 95~\textdegree C for
5~min in 1$\times$SDS buffer.

\subsection{Cell fractionation}
Fractionation followed \citet{aronica2008ema1} and earlier descriptions
(Ali et al., 2018). $1.4\times10^6$ cells were lysed in 1~mL of Lysis
buffer (50~mM Tris pH~8.0, 100~mM NaCl, 5~mM MgCl$_2$, 1~mM EDTA,
0.5\% Triton X-100, 1$\times$cOmplete protease inhibitor) on ice for
30~min and centrifuged at 15{,}000~$g$ for 10~min at 4~\textdegree C.
Proteins in 500~\textmu L of supernatant (fraction S1) were TCA-precipitated
and resuspended in 70~\textmu L 1$\times$SDS. The pellet was resuspended in
2~mL fresh Lysis buffer, sonicated ($5\times20\%$ power, 50\% pulse, five
pulses with $\geq20$~s intervals) and centrifuged at 20{,}000~$g$ for
10~min at 4~\textdegree C; 1~mL of the resulting supernatant (solubilised
chromatin, fraction S2) was TCA-precipitated and resuspended in
70~\textmu L 1$\times$SDS. For the RNase variant, cells were lysed in the
same buffer supplemented with 20~\textmu g/mL RNase~A.

\subsection{RT-PCR of pMAC-lncRNAs}
Total RNA from wild-type and \emph{EMA2} knockout cells at 6~hpm was
extracted with TRIzol. Nested PCR (95~\textdegree C/20~s,
50~\textdegree C/30~s, 68~\textdegree C/1~min; 30 cycles per round) was
performed at the M, L8 and R2 IES--MDS boundary loci using primer pairs
M5$'$-3/M3$'$-3 (M-1st), M5$'$-4/M5$'$-4 (M-2nd), L8\_5$'$-1/L8\_3$'$-1
(L8-1st), L8\_5$'$-2/L8\_3$'$-2 (L8-2nd), R2\_5$'$-1/R2\_3$'$-1 (R2-1st)
and R2\_5$'$-2/R2\_3$'$-2 (R2-2nd), resolved on a 1\% agarose gel. The
constitutively expressed \emph{RPL21} mRNA amplified with RPL21-FW and
RPL21-RV served as a loading control.

\subsection{Spt6 mutant constructs}
Surface lysines in the N, M and C regions of Spt6 were predicted with
the JASSA SUMOylation consensus tool \citep{ninova2019suvar}, and
lysine-to-arginine substitutions were introduced by overlap PCR. A
separate construct targeted every lysine within the N-terminal domain of
low complexity (DOL-KR). All constructs were selected in 300~\textmu g/mL
puromycin with 1~\textmu g/mL CdCl$_2$ induction and phenotypically
assorted up to 1.2~mg/mL puromycin.

\section{Results}

\subsection{Ema2 is expressed only during conjugation and localised
predominantly to the macronucleus}

We first asked when and where Ema2 acts. Published microarray data
\citep{miao2009microarray} indicate that \emph{EMA2} mRNA is restricted
to conjugation, so we generated an HA-tagged Ema2 strain and tracked the
protein across the life cycle. Ema2-HA was undetectable in vegetatively
growing cells; the signal first appeared in the parental MAC at 3 and
6~hpm, relocated to the new MAC at 8~hpm, and faded by 12--14~hpm
(Figure~\ref{fig:timeline}).

The timing of Ema2 localisation coincides precisely with the mid-conjugation
window in which scnRNAs enter the parental MAC and engage their genomic
targets \citep{noto2010giw1,aronica2008ema1}. Because the protein is absent
before conjugation and rapidly decays in the new MAC after 8~hpm, any
essential contribution of Ema2 to downstream development must take effect
during this narrow interval. Table~\ref{tab:tab1} summarises the expression
and localisation timeline compiled from our imaging and the published
transcriptome.

\begin{figure}[htbp]
\centering
\includegraphics[width=0.85\textwidth]{figures/fig1_ema2_timeline.pdf}
\caption{Timeline of Ema2 expression and subnuclear localisation during
conjugation. Ema2-HA is undetectable in vegetatively growing cells (Vg),
first appears in the parental MAC at 3 and 6~hpm, relocates to the new
MAC by 8~hpm, and fades by 12 and 14~hpm. Schematic compiled from the
experimental log narrative of Figures~1A--B.}
\label{fig:timeline}
\end{figure}

\begin{table}[htbp]
\centering
\caption{Ema2 expression and localisation timeline during conjugation.}
\label{tab:tab1}
\begin{tabular}{lccc}
\toprule
Stage & Time (hpm) & \emph{EMA2} mRNA & Ema2-HA localisation \\
\midrule
Vegetative & -- & Not detected & Not detected \\
Early conjugation & 3 & Induced & Parental MAC \\
Mid conjugation & 6 & Induced & Parental MAC \\
MAC transition & 8 & Induced & New MAC \\
Late conjugation & 12--14 & Declining & Faded \\
\bottomrule
\end{tabular}
\end{table}

\subsection{\emph{EMA2} is required to complete DNA elimination}

With the localisation timeline in hand, we next asked whether Ema2 is
functionally required for programmed genome rearrangement. DNA
fluorescence in situ hybridisation with Tlr1-complementary probes
\citep{cheng2010transposase} was performed on exconjugants at 36~hpm, a
time point at which wild-type cells have completed DNA elimination.
Tlr1 probe signal was confined to the MICs of wild-type exconjugants, but
in \emph{EMA2} knockout exconjugants strong signal persisted both in the
MIC and in the new MAC, indicating a failure to excise this IES family.

Quantification of FISH intensity showed that the new-MAC signal in
\emph{EMA2} knockouts was measurably lower than in \emph{TWI1} knockouts,
which completely block DNA elimination; applying the previously
established correction factor of 0.35 made the mean IES retention index
of \emph{TWI1} knockout cells equal 1, and \emph{EMA2} knockouts fell
between wild-type and \emph{TWI1} values. The \emph{EMA2} knockout
therefore represents a partial but substantial block at a step downstream
of, or parallel to, scnRNA function.

The defect in DNA rearrangement was matched by a complete reproductive
failure: no \emph{EMA2} knockout mating pair produced viable sexual
progeny in the 6-methylpurine resistance / paromomycin sensitivity assay
(Table~\ref{tab:tab2}). Together these observations establish \emph{EMA2}
as an essential factor for productive conjugation and motivated us to
ask which step of the scnRNA-directed pathway is compromised.

\subsection{Ema2 is required for target-directed small RNA degradation}

TDSD is the earliest scnRNA-dependent event in the parental MAC, so we
asked whether \emph{EMA2} loss affects it. Northern blotting with the
Mi-9 probe, which detects a population of MDS-complementary scnRNAs,
showed the expected wild-type behaviour: signal was present at 3~hpm,
reduced at 4.5~hpm and undetectable by 6~hpm. In the \emph{EMA2} knockout
cross, by contrast, the Mi-9 signal persisted through 6~hpm, matching
the qualitative schematic in Figure~\ref{fig:tdsd}.

To extend this observation genome-wide we sequenced 26- to 32-nt small
RNAs from wild-type and \emph{EMA2} knockout crosses at 3 and 8~hpm, and
counted RPKM values across tiles covering the MAC-destined sequences
(MDSs) and internal elimination sequences (IESs). MDS-complementary
scnRNAs fell sharply between 3 and 8~hpm in wild-type samples but were
only slightly reduced in \emph{EMA2} knockouts; IES-complementary
scnRNAs were comparable across strains and time points, consistent with
TDSD specifically pruning the MDS-matching pool.

The residual decrease in MDS scnRNAs that we observe in \emph{EMA2}
knockouts is milder than the complete TDSD failure seen in \emph{EMA1}
knockouts \citep{aronica2008ema1}, indicating that Ema1 and Ema2 make
distinct contributions. Because Ema1 functions as an RNA helicase that
promotes scnRNA--lncRNA base pairing, the simplest interpretation is
that Ema2 operates at a different step; the most informative candidate
is the production of the pMAC-lncRNA template itself.

\begin{figure}[htbp]
\centering
\includegraphics[width=0.75\textwidth]{figures/fig2_tdsd_schematic.pdf}
\caption{Categorical representation of Mi-9 probe scnRNA signal across
conjugation time points for wild-type and \emph{EMA2} KO cells (from the
narrative of Figure~3A). Wild-type cells show detected, reduced and absent
signal at 3, 4.5 and 6~hpm respectively, consistent with TDSD. \emph{EMA2}
KO cells retain detectable signal at all three time points.}
\label{fig:tdsd}
\end{figure}

\subsection{Ema2 is a conjugation-specific SUMO E3 ligase}

Sequence analysis of TTHERM\_00113330 revealed an SP-RING domain of the
Siz/PIAS family \citep{hochstrasser2001spring,yunus2009siz1,duan2009mms21},
but with an unusual substitution: the first zinc-binding cysteine is
replaced by histidine, and several accessory residues that stabilise the
canonical RING--E2 interface are absent. To test whether this divergent
domain retains E3 activity, we expressed recombinant GST-Ema2 and
His-Ubc9 in \emph{E.~coli} and performed a pull-down. GST-Ema2 retained
His-Ubc9 on glutathione beads whereas GST alone did not, demonstrating a
direct Ema2--Ubc9 interaction consistent with a \emph{bona fide}
SUMO~E3 \citep{sampson2001sumo,praefcke2012sumo}.

To determine whether this interaction translates into measurable E3
activity in vivo, we expressed HA-tagged Smt3 in wild-type and
\emph{EMA2} knockout strains (HA-Smt3 rescued essential \emph{SMT3}
function; Figure~EV1) and probed total SUMO conjugates by Western blot
at 4.5 and 6~hpm. High-molecular-weight (mainly $>\!200$~kDa) bands were
prominent at both time points in the wild-type cross, confirming the
conjugation-associated wave of SUMOylation described by
\citet{nasir2015sumo}. In the \emph{EMA2} knockout cross these bands
were reduced to approximately $50\%$ of the wild-type level
(Figure~\ref{fig:sumo}).

The magnitude of the defect establishes Ema2 as responsible for roughly
half of the conjugation-specific SUMO-conjugate signal, and the residual
half presumably reflects constitutive housekeeping SUMOylation and/or
other E3 activities. Because the assay captures the global SUMO
conjugate profile rather than any individual target, we next asked which
substrate(s) are specifically affected by Ema2 loss, using mass
spectrometry of His-Smt3 pulldowns.

\begin{figure}[htbp]
\centering
\includegraphics[width=0.55\textwidth]{figures/fig3_sumo_bar.pdf}
\caption{Quantification of high-molecular-weight ($>\!200$~kDa) HA-Smt3
conjugates detected in the WT cross and the \emph{EMA2} KO cross. Signals
from the WT cross were normalised to 1; the \emph{EMA2} KO cross shows
approximately 50\% of the wild-type level, as stated in the extraction of
Figure~4C.}
\label{fig:sumo}
\end{figure}

\subsection{Ema2 SUMOylates the transcription elongation factor Spt6}

Having established that Ema2 is a major source of conjugation-stage
SUMOylation, we set out to identify its substrates. His-Smt3 was
expressed and purified from wild-type and \emph{EMA2} knockout crosses
under denaturing conditions (6~M guanidine-HCl) to preserve SUMO
conjugates through Ni-NTA capture. After filtering the mass-spectrometry
hit list (excluding proteins with $\log_2$~LFQ score above~25 or more
than six consecutive histidines to remove common contaminants), the
transcription elongation factor Spt6 was the most abundant candidate
substrate in the wild-type cross and was strongly reduced in the
\emph{EMA2} knockout cross.

Direct Western analysis confirmed this result. Cells expressing HA-Spt6
in a wild-type background displayed a slower-migrating modified Spt6 band
at 4.5 and 6~hpm that was largely absent in the \emph{EMA2} knockout
cross. Anti-HA immunoprecipitation followed by anti-Smt3 blotting
identified the slower-migrating species as SUMOylated Spt6 and showed
that Ema2 accounts for the majority of Spt6 SUMOylation during
mid-conjugation. The Ema2-directed SUMOylation cascade is summarised in
Figure~\ref{fig:pathway}.

These biochemical observations place Spt6 as a principal in vivo
substrate of Ema2. Spt6 is a conserved histone chaperone and
transcription elongation factor whose SUMOylation has been reported in
other systems, so an attractive hypothesis is that modified Spt6 acts
within the parental MAC to drive pMAC-lncRNA transcription. Before
testing that hypothesis directly, we asked whether \emph{EMA2} loss
affects lncRNA output itself.

\begin{figure}[htbp]
\centering
\includegraphics[width=0.85\textwidth]{figures/fig4_pathway_schematic.pdf}
\caption{Schematic of the Ema2-directed SUMOylation cascade. Ema2
(SP-RING) binds the SUMO E2 Ubc9, which transfers Smt3 onto Spt6 and
additional substrates; SUMOylated Spt6 is retained on parental MAC
chromatin to permit pMAC-lncRNA transcription, which in turn drives TDSD
and programmed DNA elimination.}
\label{fig:pathway}
\end{figure}

\subsection{\emph{EMA2} is required for pMAC-lncRNA accumulation and
parental-MAC heterochromatin}

The hypothesis that Ema2 acts upstream of TDSD predicts that the
pMAC-lncRNA template should be absent in \emph{EMA2} knockout cells. To
test this prediction we performed nested RT-PCR across the MAC M, L8 and
R2 IES--MDS boundaries at 6~hpm. pMAC-lncRNA products were readily
detected in wild-type samples at all three loci but were undetectable
in \emph{EMA2} knockout samples, whereas the constitutively expressed
\emph{RPL21} mRNA served as a loading control and was present in both.
Genomic PCR with the same primers confirmed that the MAC loci are
present in \emph{EMA2} knockout cells (Figure~EV3), ruling out a trivial
template-loss explanation.

We extended this conclusion to total long double-stranded RNA using the
J2 antibody, which detects dsRNAs of $>40$~bp by immunofluorescence
\citep{aronica2008ema1}. Wild-type cells showed the expected wave of
nuclear dsRNA signal: MIC at 3~hpm (reflecting MIC-lncRNAs), parental
MAC at 6~hpm (pMAC-lncRNAs), and new MAC at 9~hpm (nMAC-lncRNAs). In
\emph{EMA2} knockouts the parental-MAC signal at 6~hpm was undetectable,
while the MIC signal at 3~hpm and the new-MAC signal at 9~hpm were
preserved. Importantly, \emph{EMA1} knockouts retained parental-MAC
dsRNA, indicating that Ema1 is required for scnRNA-pMAC-lncRNA
interaction rather than for pMAC-lncRNA transcription per se.

The failure to produce pMAC-lncRNAs also compromised downstream
heterochromatin formation in the parental MAC. Immunofluorescence for
H3K27me3 showed no parental-MAC signal in \emph{EMA2}, \emph{EMA1},
\emph{TWI1} or \emph{EZL1} knockout cells at 6~hpm, confirming that all
four factors converge on the same pMAC chromatin output
\citep{liu2004h3k9,liu2007rnai,frapporti2019ezl1}. This places Ema2
upstream of the Polycomb-mediated deposition of H3K27me3 in the parental
MAC, and, by extension, of the normal scnRNA--Twi1--chromatin
communication that licenses IES selection.

\begin{table}[htbp]
\centering
\caption{Phenotypes of \emph{EMA2} knockout compared with related
knockouts. ``+'' denotes the feature is present or complete, ``$-$''
denotes absent or defective, ``partial'' denotes an intermediate
phenotype, and ``n.d.'' denotes features not explicitly measured in that
strain.}
\label{tab:tab2}
\begin{tabular}{lccccc}
\toprule
Strain & Tlr1 FISH & MDS scnRNA TDSD & pMAC-lncRNA / dsRNA & pMAC H3K27me3 & Viable progeny \\
\midrule
WT & MIC only & + & + & + & + \\
\emph{EMA2} KO & partial & $-$ & $-$ & $-$ & $-$ \\
\emph{EMA1} KO & n.d. & $-$ (severe) & + & $-$ & n.d. \\
\emph{TWI1} KO & complete block & n.d. & n.d. & $-$ & n.d. \\
\emph{EZL1} KO & n.d. & n.d. & n.d. & $-$ & n.d. \\
\bottomrule
\end{tabular}
\end{table}

\subsection{Ema2 facilitates chromatin association of Spt6 and
RNA polymerase~II}

The lncRNA and SUMOylation defects together raised the possibility that
Ema2 controls a transcription-machinery step in the parental MAC. We
therefore examined the subnuclear distribution of Spt6 and of the RNAPII
subunit Rpb3 by immunofluorescence at 4.5~hpm. Spt6-HA and Rpb3 were
clearly localised to the MAC in both wild-type and \emph{EMA2} knockout
crosses (Figure~\ref{fig:phenotypes}), indicating that the defect is not
a gross mislocalisation of the transcription machinery.

Biochemical fractionation, however, revealed a selective chromatin-association
defect. $1.4\times10^6$ cells were lysed in 50~mM Tris pH~8.0, 100~mM
NaCl, 5~mM MgCl$_2$, 1~mM EDTA, 0.5\% Triton X-100 with cOmplete
protease inhibitor, yielding a soluble fraction (S1) and, after
sonication of the Triton-insoluble pellet, a solubilised chromatin
fraction (S2). In the wild-type cross, Spt6-HA and Rpb3 partitioned to
both S1 and S2, and SUMOylated Spt6 was detected exclusively in the S2
chromatin fraction. In the \emph{EMA2} knockout cross, Spt6-HA and Rpb3
in the S2 fraction were greatly reduced, while S1 levels were
essentially unchanged.

To determine whether this chromatin association requires a nascent RNA
scaffold we repeated the fractionation in the presence of 20~\textmu g/mL
RNase~A. Spt6 and RNAPII remained in the chromatin pellet under these
conditions, whereas Twi1 (which associates with chromatin via its
scnRNA--pMAC-lncRNA engagement) was released into the supernatant. The
chromatin association of Spt6 and RNAPII is therefore RNA-independent,
and Ema2 acts on a protein--chromatin interaction rather than an
RNA-tethering step. Replicate experiments supporting this conclusion
are documented in Figures~EV5 and EV6.

\begin{figure}[htbp]
\centering
\includegraphics[width=0.85\textwidth]{figures/fig5_phenotype_matrix.pdf}
\caption{Phenotype summary across knockout strains compiled from the
experimental log. A plus sign denotes the feature is present or
complete, minus denotes absent or defective, ``partial'' denotes an
intermediate phenotype (e.g.\ the DNA elimination defect of \emph{EMA2}
KO relative to \emph{TWI1} KO), and ``n.d.'' denotes features that are
not explicitly measured for that strain in the extraction.}
\label{fig:phenotypes}
\end{figure}

\subsection{Spt6 SUMOylation is dispensable for pMAC-lncRNA
transcription but contributes to DNA elimination}

The preceding experiments establish that Ema2 SUMOylates Spt6 and that
Ema2 loss collapses lncRNA transcription, but they do not yet determine
whether Spt6 SUMOylation is itself the transcriptional trigger. To
separate the two possibilities, we used the JASSA consensus algorithm
to predict surface lysines in the N-terminal, middle and C-terminal
regions of Spt6 and built lysine-to-arginine mutants in each region
(HA-SPT6-N-KR, -M-KR, -C-KR), together with a separate construct
(HA-SPT6-DOL-KR) in which every lysine in the N-terminal domain of
low complexity was substituted (Figure~\ref{fig:spt6mutants}).

All four constructs rescued vegetative \emph{SPT6} knockout lethality
(Figure~EV7), but the conjugation phenotypes differed sharply:
HA-SPT6-N-KR and HA-SPT6-M-KR strains showed severe meiotic and mating
defects, while HA-SPT6-DOL-KR and HA-SPT6-C-KR cells progressed normally
through conjugation. Crucially, HA-Spt6-DOL-KR remained SUMOylated at
wild-type levels, whereas HA-Spt6-C-KR showed no detectable SUMOylation
by anti-HA immunoprecipitation and anti-Smt3 blotting. The C-KR mutant
therefore provides a clean loss of Spt6 SUMOylation within an otherwise
functional protein.

When we examined HA-SPT6-C-KR cells, Spt6 and RNAPII remained in the
chromatin-bound S2 fraction, the parental MAC accumulated long dsRNAs
at 6~hpm, and pMAC-lncRNA production was preserved; nevertheless,
exconjugants displayed a mild DNA elimination defect. Spt6 SUMOylation
itself is therefore not required for pMAC-lncRNA transcription, yet it
still contributes to productive DNA elimination, implying that Ema2
acts through additional SUMO substrates to promote lncRNA synthesis
while Spt6 SUMOylation plays a distinct, downstream role. The full
mutant series is summarised in Table~\ref{tab:tab3}.

\begin{figure}[htbp]
\centering
\includegraphics[width=0.9\textwidth]{figures/fig6_spt6_mutants.pdf}
\caption{Summary of the Spt6 K-to-R mutant series. All constructs rescue
vegetative \emph{SPT6} KO lethality. HA-SPT6-N-KR and HA-SPT6-M-KR show
severe defects in meiotic progression and mating initiation,
respectively; HA-SPT6-DOL-KR is SUMOylated like wild-type Spt6, whereas
HA-SPT6-C-KR shows no detectable SUMOylation yet still supports
chromatin association of Spt6 and RNAPII, parental-MAC dsRNA
accumulation, and only a mild DNA elimination defect.}
\label{fig:spt6mutants}
\end{figure}

\begin{table}[htbp]
\centering
\caption{Spt6 K-to-R mutant series. Vegetative rescue refers to
\emph{SPT6} KO complementation; conjugation phenotype summarises
meiotic and mating observations; SUMOylation is scored by HA IP followed
by anti-Smt3 blot; pMAC-lncRNA and DNA elimination are scored by J2
dsRNA immunofluorescence and Tlr1 FISH, respectively.}
\label{tab:tab3}
\begin{tabular}{lcccccc}
\toprule
Mutant & Substitutions & Veg.\ rescue & Conjugation phenotype &
SUMOylation & pMAC dsRNA & DNA elimination \\
\midrule
HA-SPT6-DOL-KR & DOL lysines & + & normal & + (WT-like) & n.d. & n.d. \\
HA-SPT6-N-KR & N surface lysines & + & meiotic defect & n.d. & n.d. & n.d. \\
HA-SPT6-M-KR & M surface lysines & + & mating defect & n.d. & n.d. & n.d. \\
HA-SPT6-C-KR & C surface lysines & + & progresses & $-$ & + & mild defect \\
\bottomrule
\end{tabular}
\end{table}

\section{Discussion}

We have shown that the conjugation-specific SUMO E3 ligase Ema2 is
required for genome-wide pMAC-lncRNA transcription in
\emph{T.~thermophila}, and that this activity is a prerequisite for
scnRNA-directed chromatin communication and for programmed DNA
elimination. \emph{EMA2} knockout crosses lose MDS-targeted TDSD, fail
to accumulate pMAC-lncRNAs or parental-MAC dsRNA, lack H3K27me3 in the
parental MAC, cannot complete DNA elimination at the Tlr1 locus, and
produce no viable sexual progeny. Biochemically, Ema2 accounts for
approximately half of the high-molecular-weight ($>\!200$~kDa)
conjugation-stage SUMO conjugates, interacts directly with the SUMO E2
Ubc9, and drives the majority of Spt6 SUMOylation at 4.5 and 6~hpm.
Together these observations identify the first regulator that licenses
pMAC-lncRNA transcription and place SUMOylation upstream of small RNA
target engagement in this lineage.

The phenotypic separation between Ema1 and Ema2 clarifies how
\emph{Tetrahymena} solves the silencing paradox. Ema1 is an RNA helicase
that promotes base pairing between scnRNAs and pMAC-lncRNAs
\citep{aronica2008ema1}, so \emph{EMA1} knockouts retain parental-MAC
dsRNA but fail TDSD, whereas \emph{EMA2} knockouts lack parental-MAC
dsRNA altogether. This logical ordering makes Ema2 the earliest-acting
factor in the pMAC arm of the pathway and is consistent with the
stronger TDSD defect seen in \emph{EMA1} knockouts relative to
\emph{EMA2} knockouts. The model that emerges is one in which Ema2
supplies the lncRNA template, Ema1 mediates its base pairing with
incoming scnRNAs, and TDSD plus subsequent Polycomb-mediated H3K27me3
deposition \citep{liu2004h3k9,liu2007rnai,frapporti2019ezl1} enforce
the choice of scnRNAs that will guide IES selection in the new MAC.

Our cell-fractionation data suggest a mechanism for how SUMOylation
supports lncRNA transcription: Ema2 is required to keep Spt6 and RNAPII
associated with parental-MAC chromatin. Because the association is
preserved in the presence of RNase~A but lost in \emph{EMA2}
knockouts, Ema2 cannot merely stabilise an RNA-bridged complex; rather,
it must act on a protein--chromatin interaction during a mid-conjugation
window when the parental-MAC chromatin environment is in flux. Two
non-exclusive possibilities are that Ema2-dependent SUMOylation
counteracts a transient destabilising modification or compartmentalises
excess transcription machinery (which is upregulated during conjugation)
onto chromatin. Distinguishing these alternatives will require
time-resolved chromatin-profiling of SUMO substrates, which was beyond
the scope of the present study.

Our Spt6-C-KR mutant reveals that Spt6 SUMOylation itself is not the
direct trigger for pMAC-lncRNA transcription, because the C-KR mutant
loses detectable Spt6 SUMOylation while still retaining chromatin
association of Spt6 and RNAPII, parental-MAC dsRNA accumulation and
only a mild DNA elimination phenotype. This dissociation implies that
Ema2 must operate through additional SUMO substrates to drive lncRNA
synthesis, and that Spt6 SUMOylation plays a distinct, later role in
productive DNA elimination, perhaps acting at the nMAC stage or at an
excision-associated step. Candidate additional substrates that would
evade our His-Smt3 pulldown include SUMO conjugates that are
proteasomally degraded or that compete with other lysine modifications
for the same residue; both scenarios are consistent with the proteomic
literature on SUMO turnover \citep{praefcke2012sumo}.

In comparative terms, Ema2 represents a distinctive evolutionary
solution to the small RNA--lncRNA paradox. In fission yeast, Swi6, Chp2
and Clr4 SUMOylation stabilise heterochromatin that already templates
centromeric lncRNAs \citep{shin2005sumo}; in \emph{C.~elegans}, HDAC1
SUMOylation acts downstream of Argonaute targeting
\citep{kim2021hdac1}; and in \emph{Drosophila}, Su(var)2-10 couples
piRNA recognition to SetDB1 recruitment and H3K9me3 deposition
\citep{ninova2019suvar}, all paradigms in which SUMOylation reinforces
a pre-existing heterochromatin landscape that drives lncRNA production
\citep{lethomas2014rhino,andersen2017heterochromatin,law2013shh1,volpe2002rnai,hall2002heterochromatin}.
Ema2 differs in that it acts upstream of lncRNA transcription and does
not require a pre-formed heterochromatic template: in \emph{Tetrahymena},
H3K27me3 is established only after pMAC-lncRNAs have been transcribed and
base-paired with scnRNAs, so heterochromatin is a consequence, not a
prerequisite, of the transcription step that Ema2 licenses. The
conservation of SUMO pathway requirements for ciliate DNA elimination,
evident from RNAi depletion of \emph{UBA2} and \emph{SUMO} in
\emph{Paramecium tetraurelia} \citep{matsuda2006sumo} and from the
conjugation-stage SUMOylation wave previously characterised in
\emph{Tetrahymena} \citep{nasir2015sumo}, suggests that related E3
activities may prove to play analogous roles in other eukaryotes whose
small RNA pathways include an initial, heterochromatin-independent
licensing step.

Several limitations bound the scope of the present study. Our mass
spectrometry used stringent filters (log2 LFQ $\leq 25$ and runs of
$\leq 6$ consecutive histidines) to exclude likely contaminants, and
by design these filters may also exclude genuine Ema2 substrates that
are rapidly turned over or that contain histidine-rich stretches. Our
chromatin-association analysis reports on bulk populations of Spt6 and
Rpb3 in the parental MAC and does not yet resolve which regions of the
genome are transcribed; high-resolution RNAPII ChIP during conjugation,
and direct profiling of pMAC-lncRNA start sites, will be needed to
determine whether Ema2 licenses transcription globally or at specific
IES-associated loci. Finally, because \emph{EMA2} knockouts fail
reproductively, all of our experiments stop at the first failed
conjugation; long-term consequences for macronuclear genome integrity
over successive generations remain to be addressed in a conditional
knockdown context.

Looking forward, the demonstration that a SUMO E3 ligase licenses lncRNA
transcription opens new questions about how the parental-MAC chromatin
environment is remodelled for scnRNA engagement. Identifying the
remaining Ema2 substrates will reveal whether a single SUMO-modified
chromatin factor is responsible, or whether Ema2 modifies multiple
components of the transcription and chromatin machinery in parallel.
From a broader evolutionary perspective, small RNA-directed silencing
systems in other eukaryotes may contain analogous ``upstream''
SUMOylation steps that have so far been masked by the strong
heterochromatin-dependent phenotypes seen in later events; the Ema2
example suggests that such upstream regulators are worth a direct
search. Our work therefore provides both a specific molecular entry
point into TDSD and programmed DNA elimination in ciliates and a
general argument that the earliest licensing steps of small
RNA--chromatin communication deserve renewed attention.

\bibliography{references}

\end{document}
